\documentclass[10pt]{article}

\usepackage[letterpaper,hmargin=0.68in,vmargin=0.55in]{geometry}
\usepackage[T1]{fontenc}
\usepackage{amsmath,amssymb}
\usepackage{newtxtext,newtxmath}
\usepackage{microtype}
\usepackage{xcolor}
\usepackage[hidelinks]{hyperref}
\usepackage{enumitem}
\usepackage{titlesec}

\definecolor{DeepBlue}{HTML}{093C4C}
\definecolor{Cyan}{HTML}{087F9A}
\definecolor{SoftBlue}{HTML}{EAF5F7}
\definecolor{SoftRule}{HTML}{9CC8D3}
\definecolor{Body}{HTML}{1D3037}

\color{Body}
\setlength{\parindent}{0pt}
\setlength{\parskip}{2.5pt}
\setlength{\abovedisplayskip}{4pt}
\setlength{\belowdisplayskip}{4pt}
\setlength{\abovedisplayshortskip}{2pt}
\setlength{\belowdisplayshortskip}{2pt}
\setlist[itemize]{leftmargin=1.25em,itemsep=2pt,topsep=2pt}
\titleformat{\section}{\normalsize\bfseries\color{DeepBlue}}{}{0pt}{}
\titlespacing*{\section}{0pt}{6pt}{2pt}
\newcommand{\TV}{\operatorname{TV}}

\begin{document}

\begin{center}
  {\LARGE\bfseries\color{DeepBlue} Physical Completeness vs. Statistical Identifiability}\par
  \vspace{2pt}
  {\small A reading note on \emph{Finite Observation and Broader Physical Processes} and
  \emph{When Does a Random-Scale Flight Path Look Lognormal?}}
\end{center}

\vspace{2pt}
\noindent\fcolorbox{SoftRule}{SoftBlue}{%
\begin{minipage}{0.965\textwidth}
\small
The two papers ask different questions about what a good-looking fit means. The first asks whether a finite fit establishes that the visible physical process is the whole system; the second asks whether the distribution of observed flight trajectories is compatible with one proposed statistical population.
\end{minipage}}

\small
\section*{The distinction}

\begin{itemize}
\item \textbf{\emph{Finite Observation}: a good finite fit does not establish physical completeness.}
The example compares one visible harmonic oscillator with the same oscillator weakly coupled to a second one. During any fixed experiment, the coupling can be small enough that the two models produce nearly indistinguishable records, even though energy exchange with the second oscillator creates a large difference later:
\[
  \TV\!\left(P_\delta^{\mathcal O},P_0^{\mathcal O}\right)\leq\varepsilon,
  \qquad |q_\delta(T_*)-q_0(T_*)|=A.
\]
Here $P_\delta^{\mathcal O}$ and $P_0^{\mathcal O}$ are the probability laws of all recorded measurements under the coupled and isolated models, using the same observation protocol $\mathcal O$. The coupling-induced frequency split is $\delta>0$, $\TV$ is total-variation distance, and $\varepsilon$ is the chosen distinguishability tolerance. The corresponding visible trajectories are $q_\delta$ and $q_0$; at the later time $T_*$ they differ by the full amplitude $A$. The result does \emph{not} infer that such a larger process exists in any particular observed system. It says that finite predictive success alone cannot rule one out or certify that a physical model is complete.

The paper then states separately what must be assumed to say that further relevant scales actually exist. If $N(R,S)$ counts processes with characteristic scale $L\in(R,S]$, and their Poisson intensity is $d\eta(L)\,dL/L$, then
\[
  \Pr\{N(R,S)\geq1\}
  =1-\exp\!\left[-d\int_R^S\frac{\eta(L)}{L}\,dL\right],
  \qquad
  \int_R^\infty\frac{\eta(L)}{L}\,dL=\infty
\]
is the condition for unbounded scales almost surely. Here $R<S$ are scale cutoffs, $d$ converts scale to volume, and $\eta(L)$ is the assumed intensity per logarithmic-volume increment. This is a conditional population statement, not something established by the finite record.

\item \textbf{The aerostat paper: skewness and kurtosis reveal when one statistical population is inadequate.}
The proposed model factors each track's RMS deviation from its endpoint geodesic as
\[
  D=AQ,
  \qquad \log D=\log A+\log Q,
\]
where $D>0$ is observed RMS deviation, $A>0$ is overall route scale, and $Q>0$ is the RMS size of the dimensionless endpoint-pinned path shape. If $\log A$ is Gaussian and independent of $\log Q$, variation in $A$ should wash out the shape's non-Gaussian higher moments. With
\[
  \rho=\frac{\tau_Q^2}{\sigma_A^2+\tau_Q^2},
  \qquad \sigma_A^2=\operatorname{Var}(\log A),
  \qquad \tau_Q^2=\operatorname{Var}(\log Q),
\]
the model predicts
\[
  \operatorname{Skew}(\log D)=\operatorname{Skew}(\log Q)\rho^{3/2},
  \qquad
  \operatorname{ExKurt}(\log D)=\operatorname{ExKurt}(\log Q)\rho^2.
\]
Here $\operatorname{Var}$ is variance, $\operatorname{Skew}$ is standardized skewness, and $\operatorname{ExKurt}$ is excess kurtosis; $\rho$ is the fraction of total log-variance contributed by shape. After matching the observed variance, the single-population Brownian-bridge model predicts skewness about $0.003$ and excess kurtosis about $-0.002$. The data instead give $-2.40$ and $8.11$. This sharp mismatch statistically rejects the single-population explanation and motivates a near-geodesic boundary population or mixture. It detects missing statistical heterogeneity; it does not identify its physical cause or fully determine the mixture components.

Only after that population-level rejection does the paper ask when one flight can be assigned to the boundary component. If $C=0$ denotes that component, $\pi_0=\Pr(C=0)$ is its prior proportion, and $R_n^2=n^{-1}\sum_{j=1}^n\eta(t_j)^2$ is prefix RMS deviation after $n$ observations at times $t_j$, with cross-track deviations $\eta(t_j)$, then for a small-deviation threshold $\tau$,
\[
 \begin{aligned}
 \Pr(C=0\mid R_n\leq\tau)
 &\geq \frac{\pi_0(1-\beta_n)}
 {\pi_0(1-\beta_n)+(1-\pi_0)e^{-nI(\tau)}},\\[-1pt]
 n&\gtrsim
 \frac{\log((1-\pi_0)/\pi_0)+\log((1-\alpha)/\alpha)}{I(\tau)}.
 \end{aligned}
\]
Here $\beta_n$ bounds boundary misses, $I(\tau)>0$ is the ordinary population's exponential small-deviation rate, and $\alpha$ is the tolerated posterior error; $\gtrsim$ gives the waiting-time scale up to the $\beta_n$ correction. This follow-on result gives the time to diagnose an individual boundary track---about 41 minutes in the dataset---rather than the original evidence that a mixture is needed.
\end{itemize}

\noindent\fcolorbox{SoftRule}{SoftBlue}{%
\begin{minipage}{0.965\textwidth}
\small
\textbf{Finite Observation:} Does a good finite fit prove that no larger coupled physical process matters? No.\par
\textbf{Aerostat:} Do the observed skewness and kurtosis agree with the proposed single-population trajectory model? No---their mismatch makes additional statistical structure necessary.
\end{minipage}}

\vspace{4pt}
The papers are therefore complementary: the first warns that even an excellent finite fit cannot certify physical closure; the second shows how observable higher moments can nevertheless falsify a specific statistical explanation and reveal population heterogeneity.

\vfill
{\footnotesize\color{DeepBlue}
Companion papers: \emph{Finite Observation and Broader Physical Processes};
\emph{When Does a Random-Scale Flight Path Look Lognormal?} (J. R. Landers, 2026).}

\end{document}
